\documentclass[11pt,a4paper]{article}

% ── Codificación, fuentes y lenguaje ───────────────────────────────────────────
\usepackage[utf8]{inputenc}
\usepackage[T1]{fontenc}
\usepackage[default,scale=0.95]{sourcesanspro}
\renewcommand{\familydefault}{\sfdefault}
\usepackage[spanish,es-noshorthands,es-tabla]{babel}

% ── Geometría y espaciado ─────────────────────────────────────────────────────
\usepackage[top=2.2cm, bottom=2.5cm, left=2.2cm, right=2.2cm, headheight=15pt]{geometry}
\usepackage{setspace}
\setstretch{1.18}
\usepackage{parskip}

% ── Matemáticas y unidades ────────────────────────────────────────────────────
\usepackage{amsmath}
\usepackage{amssymb}
\usepackage{siunitx}

% ── Circuitos ─────────────────────────────────────────────────────────────────
\usepackage[european, straightvoltages]{circuitikz}

% ── Tablas ────────────────────────────────────────────────────────────────────
\usepackage{booktabs}
\usepackage{tabularx}
\usepackage{array}
\usepackage{longtable}
\usepackage{needspace}
\usepackage{colortbl}
\usepackage{enumitem}

% ── Colores, cajas e iconos ───────────────────────────────────────────────────
\usepackage[dvipsnames]{xcolor}
\usepackage{tcolorbox}
\tcbuselibrary{skins, breakable}
\usepackage{fontawesome5}
\usepackage{graphicx}

% ── Listados de código (dark-mode infográfico) ────────────────────────────────
%   Estilo base: fondo oscuro con números de línea. Los bloques lstlisting se
%   envuelven automáticamente en una caja tcolorbox redondeada.
\usepackage{listings}

% ── Secciones con barra de color (infografía) ────────────────────────────────
\usepackage{titlesec}
\titleformat{\section}
  {\normalfont\LARGE\bfseries\color{azulNoche}}
  {\colorbox{cyanNeon}{\color{fondoOscuro}\thesection}}{0.7em}{}
  [\vspace{2pt}\color{cyanNeon}\rule{\linewidth}{1.8pt}]
\titlespacing*{\section}{0pt}{24pt}{10pt}
\titleformat{\subsection}
  {\normalfont\large\bfseries\color{azulNoche}}
  {\textcolor{cyanNeon}{\thesubsection}}{0.7em}{}
  [\vspace{1pt}\color{grisLinea}\rule{\linewidth}{0.6pt}]
\titlespacing*{\subsection}{0pt}{14pt}{6pt}
\titleformat{\subsubsection}
  {\normalfont\normalsize\bfseries\color{azulNoche}}
  {\textcolor{cyanNeon}{\thesubsubsection}}{0.7em}{}
\titlespacing*{\subsubsection}{0pt}{10pt}{4pt}

% ── Cabeceras y pies de página ────────────────────────────────────────────────
\usepackage{fancyhdr}
\usepackage{lastpage}
\pagestyle{fancy}
\fancyhf{}
\renewcommand{\headrulewidth}{0pt}
\renewcommand{\footrulewidth}{0pt}
\fancyfoot[C]{%
  \begin{minipage}{\linewidth}
    \vspace{2pt}
    \color{cyanNeon}\rule{\linewidth}{1.2pt}\\[2pt]
    \centering\color{grisTexto}\footnotesize
    Página \thepage\ de \pageref{LastPage}
  \end{minipage}%
}

% ── Hiperenlaces ──────────────────────────────────────────────────────────────
\usepackage[colorlinks=true, linkcolor=azulNoche, urlcolor=cyanNeon]{hyperref}
\sloppy
\emergencystretch 3em

% ── Paleta de colores — tecnológico dark-mode ──────────────────────────────────
\definecolor{fondoOscuro}{HTML}{0D1B2A}%             Dark navy — fondo banda título
\definecolor{verdeTurquesa}{HTML}{004D40}%           Verde turquesa — bandas de marca
\definecolor{azulNoche}{HTML}{1B3A6B}%               Midnight blue — secciones, marcos
\definecolor{azulMedio}{HTML}{2A5298}%               Azul medio — variante de acento
\definecolor{cyanNeon}{HTML}{00C8E0}%                Cyan eléctrico — acentos primarios
\definecolor{verdeSignal}{HTML}{00B887}%             Verde señal — fortalezas, éxito
\definecolor{naranjaVivo}{HTML}{FF7043}%             Naranja vivo — mejoras, advertencias
\definecolor{rojoAlerta}{HTML}{E53935}%              Rojo alerta — errores
\definecolor{grisPapel}{HTML}{F4F6FA}%               Fondo tarjetas claras
\definecolor{grisLinea}{HTML}{C8D0E0}%               Bordes sutiles
\definecolor{grisTexto}{HTML}{6B7A99}%               Texto secundario
\definecolor{amarilloNota}{HTML}{FFB300}%            Notas / advertencias doradas
\definecolor{codigoFondo}{HTML}{1E2D3D}%             Fondo del bloque de código
\definecolor{codigoComentario}{HTML}{7A8FA6}%        Comentarios
\definecolor{codigoCadena}{HTML}{FF9D5C}%            Cadenas / strings
\definecolor{codigoKeyword}{HTML}{00C8E0}%           Palabras clave
\definecolor{codigoNumero}{HTML}{00E5A0}%            Números

% ── Banda de título: portada infográfica ──────────────────────────────────────
%   Diseño en 2 capas: banda de color (por defecto fondo oscuro) + franja de
%   acento cyan abajo. Uso: \bandaTitulo[color]{título}{subtítulo}.
%   El icono se puede cambiar con \renewcommand{\iconoBanda}{...} (FontAwesome).
\newcommand{\iconoBanda}{microchip}
\newcommand{\bandaTitulo}[3][fondoOscuro]{%
  \begin{tcolorbox}[
    enhanced,
    colback=#1, colframe=#1,
    boxrule=0pt, arc=8pt, outer arc=8pt,
    width=\linewidth,
    top=18pt, bottom=6pt, left=14pt, right=14pt,
    borderline south={3.5pt}{0pt}{cyanNeon},
  ]
    \begin{minipage}[c]{0.07\linewidth}
      \centering
      \textcolor{cyanNeon}{\fontsize{30}{30}\selectfont\faIcon{\iconoBanda}}
    \end{minipage}%
    \hspace{8pt}%
    \begin{minipage}[c]{0.88\linewidth}
      {\fontsize{20}{24}\selectfont\bfseries\color{white}#2}\par\vspace{5pt}%
      \textcolor{cyanNeon!80!white}{\small\faIcon{angle-right}~#3}%
    \end{minipage}
    \vspace{4pt}
  \end{tcolorbox}%
  \vspace{8pt}%
}

% ── Tarjetas de datos (infografía) ────────────────────────────────────────────
\tcbset{
  tarjetaDato/.style={
    enhanced, breakable,
    arc=6pt, outer arc=6pt,
    colback=grisPapel, colframe=azulNoche,
    boxrule=1pt,
    fonttitle=\bfseries\small, coltitle=white,
    attach boxed title to top left={yshift=-3mm, xshift=6mm},
    boxed title style={
      colback=azulNoche, colframe=azulNoche,
      arc=4pt, boxrule=0pt,
    },
    drop fuzzy shadow=azulNoche!30!white,
    top=8pt, bottom=8pt, left=10pt, right=10pt,
  },
  tarjetaIcono/.style={
    enhanced, breakable,
    arc=6pt, outer arc=6pt,
    colback=white, colframe=grisLinea, boxrule=0.8pt,
    fonttitle=\bfseries\small, coltitle=azulNoche,
    borderline west={3pt}{0pt}{cyanNeon},
    drop fuzzy shadow=azulNoche!25!white,
    top=6pt, bottom=6pt, left=10pt, right=10pt,
  },
  cajaTitulo/.style={
    enhanced, breakable,
    arc=6pt, outer arc=6pt,
    colback=azulNoche!10!grisPapel, colframe=azulNoche,
    boxrule=1pt,
    fonttitle=\bfseries, coltitle=white,
    attach boxed title to top left={yshift=-3mm, xshift=6mm},
    boxed title style={
      colback=azulNoche, colframe=azulNoche,
      arc=4pt, boxrule=0pt,
    },
    drop fuzzy shadow=azulNoche!25!white,
    top=8pt, bottom=8pt, left=10pt, right=10pt,
  },
  cajaContenido/.style={
    enhanced, breakable,
    arc=5pt, outer arc=5pt,
    colback=grisPapel, colframe=grisLinea,
    boxrule=0.8pt,
    fonttitle=\bfseries\small, coltitle=azulNoche,
    top=6pt, bottom=6pt, left=10pt, right=10pt,
  },
  cajaFortaleza/.style={
    enhanced, breakable,
    arc=6pt, outer arc=6pt,
    colback=verdeSignal!8!white, colframe=verdeSignal,
    boxrule=1pt,
    borderline west={4pt}{0pt}{verdeSignal},
    fonttitle=\bfseries\small, coltitle=verdeSignal!70!black,
    drop fuzzy shadow=verdeSignal!20!white,
    top=6pt, bottom=6pt, left=10pt, right=10pt,
  },
  cajaMejora/.style={
    enhanced, breakable,
    arc=6pt, outer arc=6pt,
    colback=naranjaVivo!8!white, colframe=naranjaVivo,
    boxrule=1pt,
    borderline west={4pt}{0pt}{naranjaVivo},
    fonttitle=\bfseries\small, coltitle=naranjaVivo!80!black,
    drop fuzzy shadow=naranjaVivo!20!white,
    top=6pt, bottom=6pt, left=10pt, right=10pt,
  },
  cajaRecomendacion/.style={
    enhanced, breakable,
    arc=6pt, outer arc=6pt,
    colback=cyanNeon!6!white, colframe=cyanNeon!60!azulNoche,
    boxrule=1pt,
    borderline west={4pt}{0pt}{cyanNeon},
    fonttitle=\bfseries\small, coltitle=azulNoche,
    drop fuzzy shadow=azulNoche!20!white,
    top=6pt, bottom=6pt, left=10pt, right=10pt,
  },
}

% ── Ícono + texto (encabezados de sección y elementos) ────────────────────────
\newcommand{\iconotexto}[2]{\textcolor{cyanNeon}{\faIcon{#1}}~#2}

% ── Listados de código: estilo dark + auto-envoltura ─────────────────────────
\lstdefinestyle{estiloCodigo}{
    backgroundcolor=\color{codigoFondo},
    basicstyle=\ttfamily\small\color{white},
    commentstyle=\color{codigoComentario}\itshape,
    stringstyle=\color{codigoCadena},
    keywordstyle=\color{codigoKeyword}\bfseries,
    numberstyle=\tiny\color{codigoComentario},
    numbers=left,
    numbersep=10pt,
    showstringspaces=false,
    breaklines=true,
    breakatwhitespace=false,
    tabsize=4,
    frame=none,
    captionpos=b,
    aboveskip=6pt,
    belowskip=4pt,
    xleftmargin=14pt,
    framexleftmargin=14pt,
    literate=
      {á}{{\'a}}1 {é}{{\'e}}1 {í}{{\'i}}1 {ó}{{\'o}}1 {ú}{{\'u}}1
      {Á}{{\'A}}1 {É}{{\'E}}1 {Í}{{\'I}}1 {Ó}{{\'O}}1 {Ú}{{\'U}}1
      {ñ}{{\~n}}1 {Ñ}{{\~N}}1 {ü}{{\"u}}1 {Ü}{{\"U}}1
      {¿}{{?`}}1 {¡}{{!`}}1
}
\lstdefinestyle{estiloPython}{
    style=estiloCodigo,
    language=Python,
}
\lstset{style=estiloCodigo}
\BeforeBeginEnvironment{lstlisting}{%
  \begin{tcolorbox}[
    enhanced, breakable,
    arc=6pt, outer arc=6pt,
    colback=codigoFondo, colframe=azulNoche,
    boxrule=1pt,
    drop fuzzy shadow=azulNoche!25!white,
    top=2pt, bottom=2pt, left=0pt, right=0pt,
  ]%
}
\AfterEndEnvironment{lstlisting}{\end{tcolorbox}}

% ── Cajas de contenido temáticas ──────────────────────────────────────────────
\newtcolorbox{cajaRecuerda}{
    enhanced, breakable,
    arc=6pt, outer arc=6pt,
    colback=azulNoche!10!grisPapel,
    colframe=azulNoche, boxrule=1pt,
    borderline west={4pt}{0pt}{cyanNeon},
    fonttitle=\bfseries\small\color{white},
    title={\faIcon{info-circle}~Recuerda},
    attach boxed title to top left={yshift=-3mm, xshift=6mm},
    boxed title style={colback=azulNoche, arc=4pt, boxrule=0pt},
    drop fuzzy shadow=azulNoche!25!white,
    left=10pt, right=10pt, top=8pt, bottom=8pt,
}

\newtcolorbox{cajaConcepto}{
    enhanced, breakable,
    arc=6pt, outer arc=6pt,
    colback=verdeSignal!8!white,
    colframe=verdeSignal, boxrule=1pt,
    borderline west={4pt}{0pt}{verdeSignal},
    fonttitle=\bfseries\small\color{white},
    title={\faIcon{lightbulb}~Concepto clave},
    attach boxed title to top left={yshift=-3mm, xshift=6mm},
    boxed title style={colback=verdeSignal!80!black, arc=4pt, boxrule=0pt},
    drop fuzzy shadow=verdeSignal!20!white,
    left=10pt, right=10pt, top=8pt, bottom=8pt,
}

\newtcolorbox{cajaEjemplo}{
    enhanced, breakable,
    arc=6pt, outer arc=6pt,
    colback=cyanNeon!5!white,
    colframe=cyanNeon!70!azulNoche, boxrule=1pt,
    borderline west={4pt}{0pt}{cyanNeon},
    fonttitle=\bfseries\small\color{fondoOscuro},
    title={\faIcon{code}~Ejemplo},
    attach boxed title to top left={yshift=-3mm, xshift=6mm},
    boxed title style={colback=cyanNeon, arc=4pt, boxrule=0pt},
    drop fuzzy shadow=azulNoche!20!white,
    left=10pt, right=10pt, top=8pt, bottom=8pt,
}

\newtcolorbox{cajaEjercicio}{
    enhanced, breakable,
    arc=6pt, outer arc=6pt,
    colback=azulMedio!7!white,
    colframe=azulMedio, boxrule=1pt,
    borderline west={4pt}{0pt}{azulMedio},
    fonttitle=\bfseries\small\color{white},
    title={\faIcon{pencil-alt}~Ejercicio},
    attach boxed title to top left={yshift=-3mm, xshift=6mm},
    boxed title style={colback=azulMedio, arc=4pt, boxrule=0pt},
    drop fuzzy shadow=azulNoche!20!white,
    left=10pt, right=10pt, top=8pt, bottom=8pt,
}

\newtcolorbox{cajaReto}{
    enhanced, breakable,
    arc=6pt, outer arc=6pt,
    colback=naranjaVivo!6!white,
    colframe=naranjaVivo, boxrule=1pt,
    borderline west={4pt}{0pt}{naranjaVivo},
    fonttitle=\bfseries\small\color{white},
    title={\faIcon{fire}~Reto},
    attach boxed title to top left={yshift=-3mm, xshift=6mm},
    boxed title style={colback=naranjaVivo, arc=4pt, boxrule=0pt},
    drop fuzzy shadow=naranjaVivo!20!white,
    left=10pt, right=10pt, top=8pt, bottom=8pt,
}

% ── Indicadores de dificultad ─────────────────────────────────────────────────
\newcommand{\facil}{\textcolor{verdeSignal}{\faIcon{circle}\,\textbf{Fácil}}}
\newcommand{\intermedio}{\textcolor{naranjaVivo}{\faIcon{circle}\,\textbf{Intermedio}}}
\newcommand{\dificil}{\textcolor{rojoAlerta}{\faIcon{circle}\,\textbf{Difícil}}}

% ─────────────────────────────────────────────────────────────────────────────
\renewcommand{\iconoBanda}{shield-alt}
\fancyhead[L]{\footnotesize\color{grisTexto}\faIcon{shield-alt}~CYBERSEGURIDAD}
\fancyhead[R]{\footnotesize\color{grisTexto}Guía del estudiante}
\hypersetup{pdftitle={Ciberseguridad Aplicada},pdfauthor={Prof. César Rodríguez}
  pdfsubject={CYBERSEGURIDAD}}

\begin{document}
\bandaTitulo[fondoOscuro]{Guía del estudiante}{Ciberseguridad Aplicada --- cómo se trabaja y se entrega}

\begin{center}\small\color{grisTexto}Prof. César Rodríguez \textbar{} ELECTRONICA \textbar{} 1 de noviembre de 2026\end{center}

\section{\iconotexto{info-circle}{Cómo funciona el curso}}

El curso se organiza en 10 semanas agrupadas en 5 fases. Cada semana tiene una práctica de laboratorio que se entrega al terminar la sesión. No hay proyecto integrador: lo que se evalúa es el trabajo de cada semana y la capacidad de sostener el criterio en la evaluación final.

\section{\iconotexto{users}{Grupos y roles}}

El curso trabaja con 4 grupos. Fija durante el semestre; cada grupo entrega sus propias prácticas.

\begin{itemize}[leftmargin=1.4em,itemsep=2pt,topsep=3pt]
  \item Rojo: Ejecuta la técnica ofensiva sobre el objetivo autorizado: prepara el entorno, lanza el ataque, registra comandos y payloads, y guarda cada salida con su hash. Aporta: Evidencia cruda en semana\_NN/evidencia y el apartado de reproducción del informe.
  \item Azul: Parte de su propia captura cruda y aplica el control equivalente: configuración, parche o regla de red. Repite el ataque para demostrar que falla y documenta el antes y el después. Aporta: La mitigación aplicada, el comando de verificación y la comparación de resultados.
  \item Púrpura: Instrumenta la detección del mismo evento: activa el registro, define la alerta y mide lo que se habría visto, incluidos los falsos positivos y negativos. Aporta: La regla o consulta de detección y la tabla de lo detectado frente a lo ejecutado.
  \item Legal y reporte: Verifica que el alcance y la autorización cubran lo ejecutado, anonimiza la evidencia, revisa la cadena de custodia e integra el informe técnico y el resumen ejecutivo. Aporta: La versión final del informe y la declaración explícita de lo que no se pudo verificar.
\end{itemize}

Los roles rotan una posición al inicio de cada fase (de la 1 a la 4), de modo que cada grupo pasa por las cuatro perspectivas. En la fase 5 no hay roles de grupo; la evaluación final es individual.

\textbf{Asignación de roles por rotación (idéntica en cada grupo)}

\begin{tabularx}{\linewidth}{@{}l X X X X@{}}
\toprule
Rotación & Integrante 1 & Integrante 2 & Integrante 3 & Integrante 4 \\
\midrule
Fase 1 & Rojo & Azul & Púrpura & Legal y reporte \\
Fase 2 & Azul & Púrpura & Legal y reporte & Rojo \\
Fase 3 & Púrpura & Legal y reporte & Rojo & Azul \\
Fase 4 & Legal y reporte & Rojo & Azul & Púrpura \\
\bottomrule
\end{tabularx}

Cada grupo crea un repositorio GitHub propio (privado, con el docente como colaborador) y lo mantiene actualizado. El docente revisa ese repositorio al finalizar cada semana de práctica.

\section{\iconotexto{code-branch}{Repositorio GitHub y entregas}}

La entrega se hace en GitHub. Un repositorio privado por grupo, con el docente como colaborador de lectura.

El docente revisa el repositorio al finalizar cada semana de práctica.

Corte de cada semana: Domingo 23:59 de la semana correspondiente; lo posterior no se corrige.

\textbf{Estructura del repositorio}

\begin{itemize}[leftmargin=1.4em,itemsep=2pt,topsep=3pt]
  \item README.md con integrantes, roles y tabla de semanas entregadas.
  \item semana\_NN/informe.md o informe.pdf con el entregable de esa práctica.
  \item semana\_NN/scripts/ con el código, con type hints y docstrings.
  \item semana\_NN/evidencia/ con salidas, capturas y hashes anonimizados.
  \item semana\_NN/README.md con el resumen de lo hecho y lo que no se pudo verificar.
\end{itemize}

\textbf{Reglas de entrega}

\begin{itemize}[leftmargin=1.4em,itemsep=2pt,topsep=3pt]
  \item Nada de secretos, tokens ni contraseñas en el repositorio; las claves van en .env ignorado.
  \item Nada de imágenes de disco, muestras de malware ni binarios: solo hashes y enlaces.
  \item Evidencia anonimizada: sin datos personales de terceros ni credenciales reales.
  \item Un commit por entrega con mensaje "semana NN: \textless{}entregable\textgreater{}".
  \item La práctica cuenta como entregada solo si está en el repositorio antes del corte.
\end{itemize}

\begin{tcolorbox}[cajaMejora]\textbf{Antes de empezar:} el trabajo se hace desde el equipo y la conexión de cada grupo. El único objetivo permitido es el laboratorio local (o el entorno propio declarado) y los sistemas con autorización escrita; nunca la infraestructura de la universidad ni de terceros.
\end{tcolorbox}

\section{\iconotexto{cube}{Puesta en marcha del laboratorio}}

El docker-compose de esta carpeta es opcional. Cada grupo puede usar su propio entorno en su PC siempre que respete el alcance autorizado.

Requiere \textasciitilde{}2 GB de RAM libre y \textasciitilde{}5 GB de disco; DVWA y Juice Shop juntos consumen menos de 700 MB.

Si se usa el laboratorio de la carpeta del curso, se levanta con docker compose -f laboratorio/docker-compose.yml up -d. Los servicios quedan accesibles solo en localhost y en la red doméstica de cada grupo.

\begin{tabularx}{\linewidth}{@{}l l r@{}}
\toprule
Servicio & Imagen & URL local\\
\midrule
objetivo-web & \texttt{vulnerables/web-dvwa:latest} & \texttt{localhost:8081} \\
objetivo-juice & \texttt{bkimminich/juice-shop:latest} & \texttt{localhost:3000} \\
objetivo-linux & \texttt{ubuntu:24.04} & \texttt{localhost:2222} \\
\bottomrule
\end{tabularx}

\begin{cajaRecuerda}Para reiniciar el laboratorio a su estado inicial: \texttt{docker compose down -v} y volver a levantarlo. Ningún cambio debe sobrevivir entre prácticas sin documentarse.
\end{cajaRecuerda}

\begin{itemize}[leftmargin=1.4em,itemsep=2pt,topsep=3pt]
  \item \{'Cada grupo trabaja desde su propio equipo y red': 'no hay red compartida del curso.'\}
  \item El laboratorio corre en localhost y nunca se expone a Internet.
  \item Solo se atacan contenedores del laboratorio o sistemas con autorización escrita y firmada.
  \item \{'Las muestras de malware no se ejecutan': 'el análisis es estático y sobre copias inertes.'\}
  \item La evidencia se anonimiza antes de entregarse; los datos personales no salen del laboratorio.
  \item Se prohíbe el uso de estas técnicas contra la infraestructura de la universidad o de terceros.
\end{itemize}

\section{\iconotexto{code}{Contrato de datos de los scripts}}

\begin{cajaConcepto}Todo script entregado imprime \textbf{un solo objeto JSON} por la salida estándar y devuelve un diccionario con estas claves. \texttt{status} vale \texttt{ok} o \texttt{error}; si es \texttt{error}, \texttt{errores} explica el motivo y el programa sale con código distinto de cero.
\begin{lstlisting}[language=Python,style=estiloPython]
{
  "status": "ok",
  "objetivo": "172.28.0.10",
  "duracion_s": 12.4,
  "hallazgos": [
    {"puerto": 80, "servicio": "http", "version": "nginx"}
  ],
  "errores": []
}
\end{lstlisting}
\end{cajaConcepto}

\begin{itemize}[leftmargin=1.4em,itemsep=2pt,topsep=3pt]
  \item Type hints en todas las funciones públicas y docstrings estilo Google.
  \item Errores manejados explícitamente: nunca se devuelve una lista vacía para ocultar un fallo.
  \item Sin rutas absolutas ni credenciales en el código.
  \item Salida determinista: mismo objetivo y mismas opciones producen la misma estructura.
\end{itemize}

\section{\iconotexto{clipboard-check}{Evaluación}}

\begin{tabularx}{\linewidth}{@{}X r@{}}
\toprule
Concepto & Peso (\%)\\
\midrule
Prácticas de laboratorio (semanas 1-9) & 35,0 \\
Controles de fase (semanas 3, 5 y 7) & 15,0 \\
Informes técnicos (semanas 5, 7 y 9) & 15,0 \\
Evaluación final práctica (semana 10) & 25,0 \\
Ética, autorizaciones y uso responsable & 10,0 \\
\bottomrule
\end{tabularx}

Cada práctica se corrige con la rúbrica de diez puntos incluida en su propio documento. Entregar fuera de plazo descuenta según el reglamento de la asignatura; entregar sin evidencia anonimizada se devuelve sin corregir.

\section{\iconotexto{balance-scale}{Ética y legalidad}}

\begin{itemize}[leftmargin=1.4em,itemsep=2pt,topsep=3pt]
  \item Cada grupo mantiene su propio repositorio GitHub y lo entrega actualizado antes del corte semanal; el docente revisa ese repositorio.
  \item \{'El trabajo se hace desde el equipo y la conexión de cada grupo': 'solo objetivos propios, el laboratorio provisto o sistemas con autorización escrita; nunca infraestructura de la universidad ni de terceros.'\}
  \item Solo se atacan objetivos propios o con autorización escrita y firmada.
  \item El laboratorio es local y nunca se expone a Internet.
  \item Las muestras de malware no se ejecutan y se manejan como copias inertes.
  \item La evidencia se anonimiza antes de entregarse.
  \item El uso indebido de las técnicas del curso tiene consecuencias académicas y legales.
  \item El código entregado usa type hints, docstrings estilo Google y manejo explícito de errores.
\end{itemize}

\section{\iconotexto{calendar-alt}{Calendario resumido}}

\begin{tabularx}{\linewidth}{@{}r p{4.1cm} X@{}}
\toprule
Sem. & Tema & Entregable\\
\midrule
1 & Panorama de la ciberseguridad y superficie de ataque & Tabla de activos y servicios, modelo STRIDE y conclusión priorizada de tres vectores. \\
2 & Criptografía aplicada y gestión de secretos & Registro de comandos, salidas y una tabla que asigne cada primitiva a la propiedad que garantiza. \\
3 & Marco legal, ético y normativo & Autorización firmada, análisis del caso y código de ética firmado por cada estudiante. \\
4 & Reconocimiento, OSINT y escaneo de puertos & Script integrado (huella digital + escáner), salida JSON, huella digital razonada, tabla comparativa con nmap y análisis de falsos positivos. \\
5 & Vulnerabilidades web según OWASP Top 10 & Informe técnico con payloads, evidencia anonimizada, causa raíz y mitigación por hallazgo. \\
6 & Hardening, firewall y detección & Auditor propio, hardening antes/después, reglas de firewall, comparativa de alcance y alertas de Suricata clasificadas. \\
7 & Desarrollo seguro y cadena de suministro & Parches, evidencia antes y después con el mismo payload, y salida de los tres controles. \\
8 & Forense digital: adquisición, artefactos y línea de tiempo & Imagen forense con hashes, formulario de custodia, línea de tiempo, relato del incidente y tabla de IoCs. \\
9 & Respuesta a incidentes y triage estático de malware & Playbook de una página, informe de triage y regla YARA probada, más post mortem. \\
10 & Evaluación final práctica de alcance cerrado & Informe técnico, resumen ejecutivo, evidencia preservada y defensa oral. \\
\bottomrule
\end{tabularx}
\end{document}
